% =====================================================================
%  Métodos de avaliação de frases em português e resultados em pipeline
%
%  Autor: Matheus Sales
%  Compila com pdfLaTeX + Biber (Overleaf: Menu > Compiler > pdfLaTeX)
%
%  Fonte do material: aplicacao_tecnicas.ipynb (52 células, 12 seções).
%  Todos os números reportados vêm da execução real do caderno.
%
%  ESTRUTURA MODULAR: cada seção vive em NN-nome.tex.
%  Para trabalhar em uma seção isolada, comente os demais \input.
% =====================================================================
\documentclass[11pt,a4paper]{article}

% ---------- codificação e idioma ----------
\usepackage[utf8]{inputenc}
\usepackage[T1]{fontenc}
\usepackage[brazil]{babel}
\usepackage{lmodern}
\usepackage{microtype}
\usepackage{csquotes}          % exigido pelo biblatex sob babel

% ---------- matemática ----------
\usepackage{amsmath,amssymb,amsthm}

% ---------- tabelas e figuras ----------
\usepackage{booktabs}
\usepackage{array}
\usepackage{multirow}
\usepackage{graphicx}
\usepackage{float}             % fornece [H]: a tabela fica ONDE foi escrita
\usepackage{xcolor}

% ---------- árvores de dependência e diagramas ----------
\usepackage{tikz}
\usepackage{tikz-dependency}
\usetikzlibrary{arrows.meta,positioning,shapes.geometric}

% ---------- código ----------
\usepackage{listings}
\usepackage[htt]{hyphenat}     % permite quebra de linha em \texttt

% ---------- geometria ----------
\usepackage{geometry}
\geometry{margin=2.5cm}

% ---------- referências cruzadas ----------
\usepackage[hidelinks]{hyperref}

% ---------- bibliografia ----------
% biblatex-abnt + biber. Para trocar de veículo, substituir style=abnt
% por style=numeric (genérico) ou pelo estilo exigido pela revista.
\usepackage[style=abnt,backend=biber]{biblatex}
\addbibresource{referencias.bib}

% =====================================================================
%  Ambientes de teorema
%  A Seção 5 (Resultados analíticos) depende deles: as proposições são
%  o resultado mais forte do artigo e não podem se dissolver em prosa.
% =====================================================================
\theoremstyle{plain}
\newtheorem{proposicao}{Proposição}
\newtheorem{corolario}[proposicao]{Corolário}

\theoremstyle{definition}
\newtheorem{definicao}{Definição}
\newtheorem{observacao}{Observação}

\renewcommand{\proofname}{Demonstração}

% =====================================================================
%  Notação — macros usadas em todo o artigo
%  Centralizadas aqui para que a notação permaneça única (o caderno
%  usava doze notações locais; o artigo usa uma).
% =====================================================================
\newcommand{\sent}{s}                                  % a frase
\newcommand{\lexico}{D}                                % o léxico
\newcommand{\raiz}{\rho}                               % raiz da árvore
\newcommand{\nucleo}[1]{h(#1)}                         % núcleo de que #1 depende
\newcommand{\tracos}[2]{\mu_{#1}(#2)}                  % traço observado
\newcommand{\tracosp}[2]{\hat\mu_{#1}(#2)}             % traço PREDITO
\newcommand{\ausente}{\bot}                            % ausência de traço
\newcommand{\SLOR}{\mathrm{SLOR}}
\newcommand{\PLL}{\mathrm{PLL}}
\newcommand{\PLLm}{\overline{\mathrm{PLL}}}
\newcommand{\surp}{\mathrm{surp}}
\newcommand{\prof}{\mathrm{prof}}
\newcommand{\orac}{\mathrm{orac}}
\newcommand{\ramif}{\mathrm{ramif}}
\newcommand{\vetor}[1]{\mathbf{v}(#1)}
\newcommand{\cosseno}{\cos}

% =====================================================================
%  Listagens de código (Apêndice A)
% =====================================================================
\lstset{
  basicstyle=\ttfamily\footnotesize,
  breaklines=true,
  columns=fullflexible,
  keepspaces=true,
  frame=single,
  framesep=4pt,
  language=Python,
  showstringspaces=false,
  extendedchars=true,
  inputencoding=utf8,          % acentos nos comentários do código (Apêndice A)
}

% Tabelas discutidas no texto imediatamente adjacente usam [H].
\newcommand{\dimtable}[1]{\small #1}

% =====================================================================
%  Capa
% =====================================================================
\title{\textbf{Métodos de avaliação de frases em português\\
       e resultados em pipeline}}

\author{Matheus Sales\thanks{Contato: matheusvmsales1@gmail.com}}

\date{}   % preencher com a data de submissão, ou manter vazio

\begin{document}
\maketitle

% ---------------------------------------------------------------------
\begin{abstract}
\noindent
Dada uma frase isolada, sem gabarito e sem referência, coloca-se a questão de quanto se
pode afirmar automaticamente a respeito de sua boa formação. Este trabalho levanta onze
métodos de avaliação aplicáveis ao português brasileiro --- léxicos, gramáticas de
padrões, analisadores sintáticos, modelos de $n$-grama, modelos mascarados e vetores
distribucionais ---, descreve cada um sob um formalismo único e submete todos a um mesmo
conjunto controlado de pares mínimos. A caracterização mostra que os métodos não são
intercambiáveis nem ordenáveis por qualidade: suas coberturas são complementares, e nem
toda medida bem definida constitui um detector de desvio. Quatro das medidas examinadas
degeneram --- uma reduz-se a função exclusiva do comprimento da frase, uma converge para
uma constante determinada por um hiperparâmetro de suavização, uma reduz-se à frequência
isolada da palavra e uma é invariante a toda permutação ---, e as quatro degenerações são
estabelecidas por demonstração, e não por medição. Dos métodos que sobrevivem à
caracterização compõe-se uma \emph{pipeline} de quatro etapas, cuja ordem decorre das
propriedades medidas e cujo veredito resulta da vacuidade de um conjunto de achados
localizados, com um único parâmetro calibrável em toda a cadeia. Reportam-se os resultados
da pipeline sobre o mesmo corpus, incluindo a propagação de erro entre etapas e os tipos
de desvio que ela não alcança.

\vspace{6pt}
\noindent\textbf{Palavras-chave:} avaliação automática de frases; pares mínimos;
aceitabilidade linguística; análise de dependências; língua portuguesa.
\end{abstract}

\vspace{12pt}

% =====================================================================
%  Corpo do artigo
%  Descomente cada \input à medida que a seção for escrita e revisada.
% =====================================================================
% =====================================================================
%  1. Introdução
% =====================================================================
\section{Introdução}
\label{sec:introducao}

A avaliação automática de uma frase isolada, sem gabarito e sem referência, constitui
problema de ordem prática antes de teórica. Um sistema que redige \emph{Há três dias o
incêndio atingiu uma área de mata nativa} a partir dos campos de um registro de eventos de
fogo precisa verificar o que produziu, e não dispõe de termo de comparação: não há tradução
de origem nem resumo humano. Métricas que pressupõem um texto de referência, como BLEU e
ROUGE, não se aplicam, e a verificação recai sobre a frase em si.

Os métodos disponíveis para essa tarefa são numerosos e heterogêneos: léxicos associados a
distância de edição~\cite{damerau1964,levenshtein1966,kukich1992}, gramáticas de padrões
escritos manualmente~\cite{naber2003,milkowski2010}, analisadores
sintáticos~\cite{nivre2016,rademaker2017}, modelos de linguagem dos quais derivam medidas de
aceitabilidade como o SLOR~\cite{pauls2012,lau2017} e a
pseudo-log-verossimilhança~\cite{wang2019,salazar2020}, e vetores
distribucionais~\cite{mikolov2013}. Não respondem, contudo, à mesma pergunta: alguns
produzem veredito exato, outros veredito estimado, e outros apenas um valor numérico.

Tendo em vista um sistema que gera textos sobre eventos de fogo em linguagem natural a
partir apenas de informações numéricas extraídas de um banco de dados, é necessária uma
metodologia parametrizada e rastreável para avaliar os textos gerados. Trata-se de um
problema de múltiplas dimensões --- gramática, estrutura, ortografia e semântica.
em cada uma dessas dimensões.

% \input{02-trabalhos-relacionados}
% \input{03-preliminares}
% =====================================================================
%  4. Metodologia — como o trabalho foi feito
% =====================================================================
\section{Metodologia}
\label{sec:metodologia}

\indent\indent O procedimento adotado parte de uma frase correta, modificada de forma controlada. Cada 
frase variante introduz um único desvio, de tipo conhecido. Todos os métodos são aplicados às
mesmas frases, e observa-se o que cada um detecta, o que deixa passar e por qual mecanismo.

A frase de referência é \emph{Há três dias o incêndio atingiu uma área de mata nativa},
redigida a partir dos campos de um registro de eventos de fogo. Dela derivam as doze
variantes apresentadas na Tabela~\ref{tab:perturbacoes}.

\begin{table}[H]
\centering
\dimtable{
\begin{tabular}{@{}llp{5.4cm}@{}}
\toprule
\textbf{Variante} & \textbf{Nível afetado} & \textbf{Alteração introduzida} \\
\midrule
\texttt{natvia}                & ortografia   & forma inexistente no léxico \\
\texttt{A três dias}           & escolha lexical & forma existente, mas inadequada \\
\texttt{uma áreas}             & concordância nominal & determinante × núcleo \\
\texttt{a área inteiro}        & concordância nominal & adjetivo posposto \\
\texttt{os incêndios atingiu}  & concordância verbal & sujeito nominal × verbo \\
\texttt{cantou}                & semântica    & restrição de seleção violada \\
\texttt{que atingiu}           & estrutura    & relativa sem oração principal \\
\texttt{Atingir\ldots}         & estrutura    & infinitivo isolado \\
ordem embaralhada              & ordem        & as onze palavras permutadas \\
\texttt{tres} + \texttt{áreas} + \texttt{natvia} & múltiplo & três desvios simultâneos \\
\midrule
\texttt{destruiu}              & ---          & paráfrase plausível (controle) \\
\texttt{O incêndio é intenso}  & ---          & predicação nominal (controle) \\
\bottomrule
\end{tabular}}
\caption{As doze variantes da frase de referência.}
\label{tab:perturbacoes}
\end{table}

As duas últimas linhas não contêm desvio e destinam-se a verificar o inverso: um método que
as acusa produz falso positivo. A forma \texttt{destruiu} constitui paráfrase igualmente
correta, e \emph{O incêndio é intenso} é predicação nominal legítima, em que o verbo não
ocupa a raiz da árvore.

Quatro frases auxiliares completam o corpus. Duas comparam comprimento de dependência com as
mesmas palavras em ordens distintas, e duas formam um par de permutação. São necessárias
porque os descritores só se distinguem quando o comprimento é mantido fixo.

Os métodos distribuem-se por quatro dimensões, cada uma associada a uma verificação
distinta. A ortografia verifica se as palavras existem; a gramática, se elas concordam entre
si; a estrutura, se a frase dispõe de predicado e como suas partes se articulam; e a
semântica, se o enunciado é plausível. A divisão não é apenas organizativa, uma vez que cada
dimensão pressupõe a anterior: não se avalia concordância sobre uma palavra que não existe,
nem plausibilidade sobre uma frase desprovida de predicado. É essa dependência que determina
a ordem das etapas da \emph{pipeline}.

Nenhum modelo foi treinado ou ajustado. Os limites observados adiante pertencem a eles, e não ao método: a substituição do analisador sintático
altera o resultado sem alterar uma linha da regra.
% =====================================================================
%  5. Os métodos — um por subseção, na ordem das quatro dimensões
%  Escopo: definição e um exemplo mínimo que a torne legível.
%  As medições sobre o corpus estão na Seção 6; a cadeia, na Seção 7.
% =====================================================================
\section{Os métodos}
\label{sec:metodos}

\indent\indent Cada subseção apresenta um método: o que ele mede, como é definido e um exemplo mínimo que
torne a definição legível. O que cada um detectou e deixou passar sobre o corpus está reunido
na Seção~\ref{sec:resultados}.

% --- ortografia ---
% =====================================================================
%  Método 1 — Léxico com distância de edição   (dimensão: ortografia)
%  Origem: células 3, 4, 5 e 6 do caderno
% =====================================================================
\subsection{Léxico com distância de edição}
\label{met:lexico}

O método não emprega modelo nem generalização. Opera sobre um conjunto finito de formas
atestadas, cada uma associada à sua frequência, e decompõe-se em dois passos encadeados.

O primeiro é a detecção. Dado o léxico $\lexico$,
%
\begin{equation}
\mathrm{erro}(w) = \begin{cases} 0, & w \in \lexico \\ 1, & w \notin \lexico \end{cases}
\end{equation}
%
O predicado é exato: a palavra consta ou não consta do conjunto. Constitui o único veredito
sem gradação nem limiar entre os onze métodos examinados.

O segundo é a correção. Seja $C_k(w)$ o conjunto das formas do léxico situadas a exatamente
$k$ edições de $w$. A busca é escalonada --- a distância 2 só é consultada quando a distância
1 não devolve candidato algum --- e, entre os candidatos admitidos, a frequência determina a
escolha:
%
\begin{equation}
\hat{w} = \arg\max_{c \,\in\, C(w)} f(c)
\end{equation}

A distância empregada é a de Damerau-Levenshtein, que computa a transposição de dois
caracteres adjacentes como uma única edição. A escolha não é detalhe de implementação: é ela
que situa \texttt{nativa} a uma edição de \texttt{natvia}. Sob métrica que não contemple a
transposição, a permuta custaria 2, \texttt{nativa} ficaria fora do conjunto de candidatos e
a correção emitida seria \texttt{naevia}.

O algoritmo não busca formas semelhantes. Ele enumera as 949 sequências situadas a uma
edição de \texttt{natvia} e submete cada uma ao teste de pertinência a $\lexico$; quatro
sobrevivem, conforme a Tabela~\ref{tab:candidatos}.

\begin{table}[H]
\centering
\dimtable{
\begin{tabular}{@{}lll r@{}}
\toprule
$c$ & Operação sobre \texttt{natvia} & & $f(c)$ \\
\midrule
\texttt{nativa} & transposição de \texttt{v} e \texttt{i} & & 912 \\
\texttt{naevia} & substituição de \texttt{t} por \texttt{e}  & & 103 \\
\texttt{natia}  & remoção do \texttt{v}              & &  77 \\
\texttt{latvia} & substituição de \texttt{n} por \texttt{l}  & &  23 \\
\bottomrule
\end{tabular}}
\caption{Candidatos a uma edição de \texttt{natvia}, com sua frequência no léxico.}
\label{tab:candidatos}
\end{table}

Os quatro candidatos equidistam da forma detectada, de modo que a decisão recai integralmente
sobre $f$. A frequência não opera como critério auxiliar de desempate, mas como único
critério.


% --- gramática ---
% =====================================================================
%  Método 2 — Traços morfossintáticos sobre a árvore  (dimensão: gramática)
%  Escopo: o que a técnica se propõe a fazer e as fórmulas.
%  A aplicação das fórmulas está na subseção de resultados.
% =====================================================================
\subsection{Traços morfossintáticos sobre a árvore de dependências}
\label{met:tracos}

\indent\indent O método se propõe a verificar concordância. A concordância não constitui heurística passível
de aproximação: define-se como igualdade de traços entre dois nós ligados na árvore sintática,
e o método transcreve literalmente essa definição.

O analisador~\cite{nivre2016} associa a cada palavra $t$ três informações: a classe gramatical,
o núcleo $\nucleo{t}$ do qual ela depende e um mapa de traços
%
\begin{equation}
\tracosp{\tau}{t} \in \{\text{valor},\ \ausente\}
\label{eq:tracos-mapa}
\end{equation}
%
em que $\tau$ designa o traço --- \texttt{Number}, \texttt{Gender} ou \texttt{Person} --- e
$\ausente$ indica ausência, uma vez que preposições e advérbios, entre outras classes, não
portam número nem gênero. A árvore é o conjunto de arcos $A = \{(t,\, r,\, \nucleo{t})\}$,
sendo $r$ a função sintática. Emprega-se $\tracosp{\tau}{t}$, e não $\tracos{\tau}{t}$, porque
o traço não é observado, mas \emph{predito} pelo anotador --- distinção notacional aqui e
consequente na Seção~\ref{sec:discussao}, onde um traço mal predito suprime um desvio real.

Nem todo arco impõe concordância, e cada tipo impõe traços distintos. Seja $T(r)$ o conjunto
de traços cuja coincidência o arco $r$ exige:
%
\begin{equation}
T(r) = \begin{cases}
\{\texttt{Number},\ \texttt{Gender}\}, & r \in \{\texttt{det},\ \texttt{amod},\ \texttt{acl}\},\ \nucleo{t} \text{ é NOUN/PROPN} \\[3pt]
\{\texttt{Number},\ \texttt{Person}\}, & r \in \{\texttt{nsubj},\ \texttt{nsubj:pass}\},\ \nucleo{t} \text{ é VERB/AUX} \\[3pt]
\varnothing, & \text{caso contrário}
\end{cases}
\label{eq:tracos-Tr}
\end{equation}
%
O conjunto de erros da frase é
%
\begin{equation}
E(\sent) = \bigl\{\, (t,\,h,\,\tau) \;:\; (t,r,h) \in A,\ \ \tau \in T(r),\ \
\tracosp{\tau}{t} \neq \ausente,\ \ \tracosp{\tau}{h} \neq \ausente,\ \
\tracosp{\tau}{t} \neq \tracosp{\tau}{h} \,\bigr\}
\label{eq:tracos-E}
\end{equation}
%
e a frase é aceita quando $E(\sent) = \varnothing$. A saída não é um escore, mas o conjunto das
palavras responsáveis, cada uma acompanhada do termo com que diverge e do traço violado.

Três decisões estão inscritas na formulação. A guarda $\tracosp{\tau}{\cdot} \neq \ausente$ é
obrigatória: sem ela, toda palavra desprovida de número seria computada como discordante de
seu núcleo. O terceiro caso, $T(r) = \varnothing$, previne o falso positivo mais frequente ---
em \emph{as flores do jardim}, \texttt{jardim} liga-se a \texttt{flores} por \texttt{nmod}, e
substantivos assim ligados não concordam entre si. E a inclusão de \texttt{acl} não é
acessória: em \emph{os dados coletados} o particípio recebe essa função e a classe VERB, e não
\texttt{amod}, de modo que sua omissão suprimiria a verificação de concordância de particípio.

% =====================================================================
%  Método 3 — Regras escritas à mão (LanguageTool)  (dimensão: gramática)
%  Escopo: o que a técnica se propõe a fazer e as fórmulas.
%  A aplicação das fórmulas está na subseção de resultados.
% =====================================================================
\subsection{Regras escritas à mão}
\label{met:regras}

\indent\indent O método se propõe a confrontar a frase com a norma escrita. A gramática normativa
encontra-se codificada em compêndios e gramáticas de referência, e o método não realiza
inferência: cada regra é transcrita manualmente como um padrão sobre a sequência anotada, e o
verificador busca casamentos~\cite{naber2003,milkowski2010}.

A frase é tokenizada e anotada, constituindo uma sequência de triplas
%
\begin{equation}
\sent = \bigl\langle (w_1, \ell_1, p_1),\ (w_2, \ell_2, p_2),\ \ldots,\ (w_n, \ell_n, p_n) \bigr\rangle
\label{eq:regras-seq}
\end{equation}
%
em que $w_i$ designa a forma escrita, $\ell_i$ o lema e $p_i$ a etiqueta morfológica.

Cada regra $r$ consiste em um padrão $\pi_r$: expressão sobre essas triplas, análoga a uma
expressão regular, porém operando sobre atributos linguísticos em lugar de caracteres. A
regra dispara nas posições em que o padrão casa,
%
\begin{equation}
D_r(\sent) = \{\, i \;:\; \pi_r \text{ casa em } \sent \text{ a partir da posição } i \,\}
\label{eq:regras-Dr}
\end{equation}
%
e a saída do método é a união sobre a gramática $R$, o conjunto finito de regras elaboradas
para o idioma:
%
\begin{equation}
D(\sent) = \bigcup_{r \in R} \bigl\{ (i,\, r) \;:\; i \in D_r(\sent) \bigr\}
\label{eq:regras-D}
\end{equation}

A propriedade determinante do método é que $R$ é finito e de elaboração manual. Daí decorre,
sem gradação intermediária,
%
\begin{equation}
\text{erro} \notin \bigcup_{r \in R} \mathcal{L}(\pi_r) \;\Longrightarrow\; \text{invisível}
\label{eq:regras-invisivel}
\end{equation}
%
em que $\mathcal{L}(\pi_r)$ designa o conjunto de construções que o padrão reconhece. Não há
generalização nem grau de confiança: na ausência de regra correspondente, o desvio não é
detectado.

A cobertura do método equivale, portanto, ao trabalho humano investido no idioma.

% =====================================================================
%  Método 4 — SLOR   (dimensão: gramática)
%  Escopo: o que a técnica se propõe a fazer e as fórmulas.
%  A aplicação das fórmulas está na subseção de resultados.
% =====================================================================
\subsection{SLOR — Syntactic Log-Odds Ratio (razão sintática de log-chances)}
\label{met:slor}

\indent\indent O método se propõe a separar \emph{raro} de \emph{incorreto}. Uma frase malformada é
improvável sob um modelo de linguagem~\cite{pauls2012,lau2017}, o que sugeriria a leitura
direta de $\ln P(\sent)$; a probabilidade bruta, porém, não faz essa distinção, uma vez que
uma frase correta de vocabulário técnico é improvável por conter palavras improváveis, e não
por má construção. O SLOR desconta a contribuição atribuível à frequência isolada das palavras
e normaliza pelo comprimento:
%
\begin{equation}
\SLOR(\sent) = \frac{\ln P_{\mathrm{LM}}(\sent) \;-\; \ln P_{\mathrm{uni}}(\sent)}{|\sent|}
\label{eq:slor}
\end{equation}

O numerador confronta dois termos. O primeiro corresponde ao modelo sensível ao contexto, aqui
um bigrama; o segundo, ao mesmo texto sob modelo que ignora a ordem:
%
\begin{equation}
\ln P_{\mathrm{LM}}(\sent) = \sum_{i=1}^{|\sent|} \ln P(w_i \mid w_{i-1})
\qquad\qquad
\ln P_{\mathrm{uni}}(\sent) = \sum_{i=1}^{|\sent|} \ln P(w_i)
\label{eq:slor-termos}
\end{equation}
%
A subtração isola a contribuição da estrutura. Sob um bigrama exato, a coincidência dos dois
termos anularia o SLOR, e o sinal seria interpretável: positivo indicaria que a ordem
acrescenta probabilidade, negativo indicaria sequência anti-estrutural.

As estimativas derivam de contagens em corpus, sem recurso a modelo treinado:
%
\begin{equation}
P(w) = \frac{c(w) + 1}{N + V}
\qquad\qquad
P(w_i \mid w_{i-1}) = \lambda\,\frac{c(w_{i-1}, w_i)}{c(w_{i-1})} \;+\; (1-\lambda)\,P(w_i)
\label{eq:slor-estimadores}
\end{equation}
%
Os termos $+1$ e $+V$ constituem a suavização de Laplace, que impede que uma palavra nunca
observada anule a estimativa. A interpolação por $\lambda$ resolve o problema análogo no
bigrama: quando o par $(w_{i-1}, w_i)$ não foi observado, a estimativa recua para a
distribuição unigrama. Adotou-se $\lambda = 0{,}7$ sobre o mac\_morpho, com $N = 1.012.661$
palavras, $V = 53.115$ formas distintas e 51.115 frases.

Desse recuo decorre uma propriedade do estimador que a definição original não possui. Quando
$c(w_{i-1}, w_i) = 0$, resta $P(w_i \mid w_{i-1}) = (1-\lambda)P(w_i)$, e a contribuição da
posição ao numerador vale
%
\begin{equation}
\ln P(w_i \mid w_{i-1}) - \ln P(w_i) = \ln(1-\lambda) = -1{,}204
\label{eq:piso-slor}
\end{equation}
%
O piso do SLOR não é, portanto, zero, mas $\ln(1-\lambda)$: uma frase cujos bigramas jamais
foram observados recebe $-1{,}204$ independentemente de estar bem construída. O método devolve
um número por frase, sem indicação de posição.

% =====================================================================
%  Método 5 — PLL   (dimensão: gramática)
%  Origem: células 20, 22 e 23 do caderno
% =====================================================================
\subsection{PLL --- pseudo-log-verossimilhança}
\label{met:pll}

O SLOR estima a probabilidade sob um modelo causal, restrito ao contexto à esquerda. Um
modelo mascarado dispõe de ambos os lados: suprime-se uma palavra e consulta-se o modelo
sobre o valor esperado naquela posição, com a totalidade da frase remanescente disponível. A
pseudo-log-verossimilhança soma essas respostas, uma por posição:
%
\begin{equation}
\PLL(\sent) = \sum_{i=1}^{|\sent|} \ln P_{\mathrm{MLM}}\bigl(\sent_i \mid \sent_{\setminus i}\bigr)
\end{equation}
%
em que $\sent_{\setminus i}$ designa a frase com a posição $i$ substituída pela marca
\texttt{[MASK]}. Cada parcela exige uma passagem pelo modelo, de modo que uma frase de $n$
palavras demanda $n$ inferências. Empregou-se o BERTimbau \emph{base cased}~\cite{souza2020}.

O qualificativo \emph{pseudo} é literal, e não ressalva notacional. As $|\sent|$ parcelas se
sobrepõem, pois cada uma pressupõe corretas todas as demais posições. A soma, portanto, não
corresponde a $\ln P(\sent)$: não há regra de cadeia que a produza, e ela não integra 1 sobre
o espaço das frases. Decorre daí que o valor absoluto carece de significado, e apenas
comparações são interpretáveis.

A soma é grandeza extensiva, pois cada posição acrescenta uma parcela negativa e o escore
decresce com o comprimento ainda que a frase seja bem formada; a comparação de frases de
extensões distintas pela soma mensura sobretudo a extensão. Adota-se, por isso, a média, que
é intensiva e torna a medida homogênea ao SLOR, igualmente dividido por $|\sent|$:
%
\begin{equation}
\PLLm(\sent) = \frac{1}{|\sent|}\sum_{i=1}^{|\sent|} \ln P_{\mathrm{MLM}}\bigl(\sent_i \mid \sent_{\setminus i}\bigr)
\end{equation}

Confrontando as duas medidas sob a mesma forma, evidencia-se a diferença que determina o
emprego de cada uma:
%
\begin{equation}
\SLOR(\sent) = \frac{\ln P_{\mathrm{LM}}(\sent) - \overbrace{\ln P_{\mathrm{uni}}(\sent)}^{\text{referência}}}{|\sent|}
\qquad\qquad
\PLLm(\sent) = \frac{\ln P_{\mathrm{MLM}}(\sent) - \overbrace{0}^{\text{nenhuma}}}{|\sent|}
\end{equation}
%
O PLL não subtrai referência alguma. Trata-se da mesma forma com a linha de base anulada, e
essa ausência constitui simultaneamente sua vantagem --- nada há a estimar a partir de
contagens --- e sua limitação.

O método devolve um número por frase. A Tabela~\ref{tab:pll} apresenta o resultado sobre o
par mínimo de concordância verbal.

\begin{table}[H]
\centering
\dimtable{
\begin{tabular}{@{}lrr@{}}
\toprule
\textbf{Frase} & \textbf{Soma} & $\PLLm$ \\
\midrule
Há três dias o incêndio atingiu uma área de mata nativa.   & $-9{,}16$  & $-0{,}763$ \\
Há três dias os incêndios atingiu uma área de mata nativa. & $-15{,}16$ & $-1{,}263$ \\
\bottomrule
\end{tabular}}
\caption{O PLL sobre o par mínimo de concordância verbal ($|\sent| = 12$ em ambas).}
\label{tab:pll}
\end{table}

A separação é nítida, e a queda concentra-se em uma única posição: a palavra \texttt{atingiu}
passa de $-0{,}91$ na versão correta para $-6{,}09$ na versão com desvio. Manifesta-se aqui a
bidirecionalidade do modelo, uma vez que, ao mascarar o verbo, ele dispõe de \emph{os
incêndios} à esquerda e conclui pela inadequação do singular.


% --- estrutura ---
% =====================================================================
%  Método 6 — Teste do núcleo predicativo   (dimensão: estrutura)
%  Origem: células 25, 26 e 27 do caderno
% =====================================================================
\subsection{Teste do núcleo predicativo}
\label{met:raiz}

Toda frase requer um predicado. Não se trata de heurística nem de tendência estatística, mas
de propriedade sintática, e em uma árvore de dependências essa propriedade é verificável em
um único ponto: a raiz, que por construção é o nó do qual todos os demais dependem.

Seja $\raiz(\sent)$ a raiz da árvore e $\mathrm{filhos}(\raiz)$ seus dependentes diretos. A
frase possui predicado quando
%
\begin{equation}
\mathrm{Pred}(\sent) \;=\;
\underbrace{\bigl[\mathrm{pos}(\raiz) \in \{\texttt{VERB}, \texttt{AUX}\}\bigr]}_{\text{predicado verbal}}
\;\lor\;
\underbrace{\bigl[\exists\, c \in \mathrm{filhos}(\raiz) : \mathrm{dep}(c) = \texttt{cop}\bigr]}_{\text{predicado nominal}}
\end{equation}
%
Duas linhas de código e um único nó inspecionado bastam, o que faz deste o método de menor
custo entre os onze e ainda assim produz veredito binário e localizado.

A inspeção recai sobre a raiz, e não sobre a presença de verbo, porque o teste ingênuo ---
verificar se existe algum verbo na frase --- aprova o fragmento. Confrontem-se \emph{O
incêndio \textbf{atingiu} a mata nativa} e \emph{O incêndio que \textbf{atingiu} a mata
nativa}: ambas contêm o verbo \texttt{atingiu}, mas a primeira constitui oração e a segunda
não. O que difere não é a presença do verbo, e sim sua posição na árvore, pois o pronome
\texttt{que} o rebaixa a oração relativa (\texttt{acl:relcl}) e promove o substantivo à raiz.
Uma única palavra converte oração em fragmento, e apenas a estrutura registra o fato.

O segundo termo da disjunção é igualmente necessário. Em predicações nominais, como \emph{O
carro é azul}, a raiz é o predicativo \texttt{azul}, etiquetado \texttt{ADJ}, e o verbo
\texttt{é} figura como dependente com função \texttt{cop}. Na ausência dessa cláusula, toda
construção copular seria classificada como fragmento, o que inviabilizaria o método em
português.

O método devolve um veredito binário acompanhado do nó que o justifica, conforme a
Tabela~\ref{tab:raiz}.

\begin{table}[H]
\centering
\dimtable{
\begin{tabular}{@{}llll@{}}
\toprule
\textbf{Frase} & \textbf{Raiz} & \textbf{Classe} & \textbf{Veredito} \\
\midrule
Há três dias o incêndio atingiu uma área\ldots       & \texttt{atingiu} & VERB & ORAÇÃO \\
Há três dias o incêndio que atingiu uma área\ldots   & \texttt{dias}    & NOUN & \textbf{FRAGMENTO} \\
O incêndio é intenso.                                & \texttt{intenso} & ADJ + cópula & ORAÇÃO \\
\bottomrule
\end{tabular}}
\caption{O teste da raiz sobre três frases do corpus.}
\label{tab:raiz}
\end{table}

A segunda linha corresponde ao fragmento, e a terceira exercita a cláusula da cópula: a raiz é
um adjetivo, o primeiro termo da disjunção é falso, e ainda assim a frase é corretamente
aprovada. O método mensura completude, e nada além disso, não estabelecendo nada quanto à
correção ou à plausibilidade da frase.

% =====================================================================
%  Método 8 — Comprimento de dependência e não-projetividade  (estrutura)
%  Origem: células 33, 34 e 35 do caderno
% =====================================================================
\subsection{Comprimento de dependência e não-projetividade}
\label{met:dependencia}

Os métodos anteriores examinavam quais palavras se vinculam a quais. Este examina a distância
que as separa na ordem linear e a eventual sobreposição entre os vínculos. São duas medidas
extraídas da mesma árvore, ambas operando sobre apenas duas informações --- o conjunto de
arcos e a posição de cada palavra ---, sem que intervenham etiquetas, morfologia ou léxico.

A primeira é o comprimento de dependência. Cada palavra $t$, exceto a raiz, vincula-se ao seu
núcleo $\nucleo{t}$; sendo $i(\cdot)$ a posição na frase, o arco apresenta comprimento
$d(t) = \lvert\, i(t) - i(\nucleo{t}) \,\rvert$, e a soma sobre a frase integral fornece
%
\begin{equation}
D(\sent) = \sum_{t \,\neq\, \raiz} d(t)
\qquad\qquad
\bar{D}(\sent) = \frac{D(\sent)}{|A|}
\end{equation}
%
em que $|A|$ designa o número de arcos, excluída a pontuação. Uma vez que $D$ cresce com a
extensão da frase, apenas $\bar{D}$ é comparável entre frases de comprimentos diferentes.

A justificativa da medida é de ordem cognitiva. Um arco vincula duas palavras cuja relação é
necessária à interpretação; ao processar a primeira, o leitor mantém o vínculo em aberto até
encontrar a segunda, e quanto maior o arco, maior o intervalo de retenção. É a hipótese da
minimização do comprimento de dependência~\cite{gibson1998,liu2008}, para a qual há evidência
em larga escala em dezenas de línguas~\cite{futrell2015}.

A segunda medida é a não-projetividade. Tomem-se dois arcos, cada um representado com a
extremidade de menor índice em primeiro lugar, $(i,j)$ e $(k,l)$, com $i<j$ e $k<l$. Os arcos
se cruzam quando exatamente uma extremidade de cada um recai no interior do outro:
%
\begin{equation}
i < k < j < l \qquad\text{ou}\qquad k < i < l < j
\end{equation}
%
Sendo $X(\sent)$ o total de pares que se cruzam, a árvore é projetiva quando $X(\sent) = 0$,
isto é, quando é possível traçar todos os arcos acima da linha sem sobreposição entre eles. A
condição é puramente aritmética sobre índices, não comportando elemento linguístico algum.

O corpus contém duas frases com as mesmas 13 palavras e os mesmos 12 arcos, ambas corretas,
diferindo apenas na ordem. A Tabela~\ref{tab:dependencia} apresenta as duas medidas sobre
elas.

\begin{table}[H]
\centering
\dimtable{
\begin{tabular}{@{}lrrrr@{}}
\toprule
\textbf{Frase} & $D$ & $\bar{D}$ & \textbf{Maior arco} & $X$ \\
\midrule
O vento agravou a seca que intensificou o fogo que destruiu a mata. & 17 & 1,42 & 2 & 0 \\
O fogo que a seca que o vento agravou intensificou destruiu a mata. & 43 & 3,58 & 9 & 0 \\
\bottomrule
\end{tabular}}
\caption{As duas medidas sobre o mesmo conteúdo em ordens distintas.}
\label{tab:dependencia}
\end{table}

Na primeira frase nenhum arco excede 2, pois cada relativa se vincula ao substantivo
imediatamente anterior e a frase se resolve da esquerda para a direita sem retenção. Na
segunda os vínculos se acumulam --- \texttt{fogo} encontra seu verbo nove posições adiante
--- e $\bar{D}$ mais que duplica. Como referência aproximada, a prosa comum em português
situa-se entre 1 e 2; acima de 3, a frase mantém número elevado de vínculos simultaneamente
em aberto.

Este é o aspecto inacessível aos índices da subseção anterior. Ambas as frases encaixam duas
relativas e receberiam a mesma profundidade e o mesmo número de orações, de modo que o que as
distingue é a posição em que as relativas foram inseridas. Já $X$ é nulo em ambas, como
ocorre na quase totalidade dos casos, o que a Seção~\ref{sec:empiricos} retoma.


% --- semântica ---
% =====================================================================
%  Método 9 — Surpresa por token   (dimensão: semântica)
%  Origem: células 41, 42 e 43 do caderno
% =====================================================================
\subsection{Surpresa por token}
\label{met:surpresa}

O SLOR e o PLL devolvem um valor para a frase integral. Prestam-se à ordenação de frases, mas
não localizam o desvio. A surpresa devolve um valor por palavra: é a mesma grandeza,
examinada antes da agregação.

Para cada posição, a surpresa é o logaritmo negativo da probabilidade que o modelo atribui à
palavra efetivamente ocorrente:
%
\begin{equation}
\surp(w_i) = -\ln P\bigl(w_i \mid w_{i-1}\bigr)
\end{equation}
%
Como $P \in (0,1]$, o logaritmo é negativo e a surpresa é sempre não negativa. Assume valor
nulo sob certeza do modelo e cresce ilimitadamente à medida que a palavra se torna improvável.
A unidade é relevante: em logaritmo natural a grandeza é expressa em \emph{nats} e, sob
$\log_2$, em \emph{bits}, caso em que a leitura é literal --- a quantidade de informação que
a palavra veicula. A surpresa corresponde exatamente à autoinformação da teoria da informação.

O método não constitui técnica nova, mas a parcela que o SLOR já agregava. Retomando os dois
estimadores da subseção~\ref{met:slor},
%
\begin{equation}
\SLOR(\sent) = \frac{1}{|\sent|}\sum_{i} \Bigl[\,\surp_{\mathrm{uni}}(w_i)
- \surp_{\mathrm{bi}}(w_i)\,\Bigr]
\end{equation}
%
isto é, o SLOR é a média da redução de surpresa produzida pelo contexto. Este método é o mesmo
cálculo exposto antes da agregação, de modo que nada é recomputado e nenhum modelo adicional
se faz necessário.

A leitura proposta consiste em identificar o pico,
%
\begin{equation}
\hat{\imath} \;=\; \arg\max_{i} \; \surp(w_i)
\end{equation}
%
sob a expectativa de que a surpresa se eleve na posição defeituosa, o que permitiria localizar
a palavra. Convém registrar a suposição embutida nessa leitura: assume-se que \emph{improvável}
e \emph{incorreto} coincidam.

O método devolve um perfil com um valor por palavra, do qual se retém o pico. A
Tabela~\ref{tab:surpresa} apresenta o resultado sobre três frases.

\begin{table}[H]
\centering
\dimtable{
\begin{tabular}{@{}lrl@{}}
\toprule
\textbf{Frase} & \textbf{Pico} & \textbf{Palavra} \\
\midrule
Há três dias o incêndio atingiu uma área de mata nativa & 11,47 & \texttt{atingiu} \\
Há três dias o incêndio cantou uma área de mata nativa  & 13,70 & \texttt{cantou} \\
Focos recorrentes apresentaram persistência anômala     & 15,08 & \texttt{anômala} \\
\bottomrule
\end{tabular}}
\caption{O pico de surpresa em três frases, em nats.}
\label{tab:surpresa}
\end{table}

Nas posições em que o contexto efetivamente intervém, a leitura é adequada. Isoladamente,
\texttt{nativa} é palavra de frequência muito baixa, com surpresa de $11{,}58$; na sequência
de \texttt{mata}, contudo, é praticamente esperada e o valor decresce para $3{,}53$. É o
comportamento pretendido pela técnica, a saber, mensurar a palavra em posição, e não em
abstrato. As duas outras linhas da tabela, entretanto, correspondem a frases corretas, e a
Seção~\ref{sec:empiricos} examina o que o $\arg\max$ está de fato selecionando.

% =====================================================================
%  Método 10 — Plausibilidade do verbo mascarado   (dimensão: semântica)
%  Origem: células 45, 46 e 47 do caderno
% =====================================================================
\subsection{Plausibilidade do verbo mascarado}
\label{met:verbo}

Os métodos anteriores atribuem escore ao que está escrito. Este inverte o procedimento:
suprime-se a palavra e determina-se qual palavra seria esperada naquela posição, confrontando
o texto com a expectativa apenas em seguida.

A supressão de posição arbitrária não seria adequada, e a escolha é sintática: suprime-se a
raiz da árvore de dependências, isto é, o verbo que sustenta a predicação e seleciona os
demais argumentos. O método só se aplica quando $\mathrm{pos}(\raiz(\sent)) = \texttt{VERB}$
e, quando essa condição é falsa --- fragmento ou predicação nominal ---, a resposta adequada é
\texttt{n/a}, e não um escore. Não se trata de limitação implícita, mas de elemento
constitutivo da definição. Cabe observar o arranjo: o analisador sintático determina a posição
a ser consultada, e o modelo estatístico determina o que é plausível naquela posição, sendo
que nenhum dos dois executaria a tarefa isoladamente.

Substituída a raiz pela marca, obtém-se $\sent_{\setminus\raiz}$, e o modelo devolve uma
distribuição sobre todo o vocabulário $\mathcal{V}$, que soma 1. Diferentemente do PLL,
trata-se de probabilidade em sentido estrito, uma vez que a soma percorre as alternativas de
uma única posição, que é exatamente o que o modelo mascarado foi treinado para prever. Dela se
extraem duas leituras, o logaritmo da probabilidade do verbo escrito e a margem até o melhor
preenchimento possível:
%
\begin{equation}
\hat{v} = \arg\max_{v \in \mathcal{V}} P\bigl(v \mid \sent_{\setminus\raiz}\bigr)
\qquad\qquad
m(\sent) = \ln P\bigl(\hat{v} \mid \sent_{\setminus\raiz}\bigr) - \ln P\bigl(v_\raiz \mid \sent_{\setminus\raiz}\bigr)
\end{equation}
%
Por construção $m(\sent) \geq 0$, e $m(\sent) = 0$ exatamente quando o verbo escrito é a
primeira escolha do modelo.

Adota-se a margem, e não o valor absoluto, porque este depende do grau geral de
previsibilidade da posição: em contexto muito restrito toda alternativa recebe probabilidade
elevada e, em contexto vago, toda alternativa recebe probabilidade baixa, inclusive o verbo
adequado. A margem normaliza pelo máximo daquela posição, quantificando a distância ao melhor
preenchimento em lugar da probabilidade absoluta.

O método devolve a margem, acompanhada da lista das formas esperadas, conforme a
Tabela~\ref{tab:verbo}.

\begin{table}[H]
\centering
\dimtable{
\begin{tabular}{@{}lrr@{}}
\toprule
\textbf{Frase} & $\ln P$ (escrito) & \textbf{Margem} \\
\midrule
Há três dias o incêndio \textbf{destruiu} uma área de mata nativa. & $-0{,}72$  & $0{,}00$ \\
Há três dias o incêndio \textbf{atingiu} uma área de mata nativa.  & $-0{,}91$  & $0{,}20$ \\
Há três dias o incêndio \textbf{cantou} uma área de mata nativa.   & $-14{,}45$ & $\mathbf{13{,}73}$ \\
\bottomrule
\end{tabular}}
\caption{A margem do verbo da raiz. As formas esperadas são as mesmas nas três frases:
\texttt{[destruiu, atingiu, atinge, invadiu, em]}.}
\label{tab:verbo}
\end{table}

A separação é ampla, com quase 14 nats entre o verbo adequado e o anômalo e os dois verbos
plausíveis situados junto ao zero. O aspecto metodologicamente relevante é a lista de formas
esperadas ser idêntica nas três frases, o que não poderia ser diferente: ao mascarar a raiz,
suprime-se justamente a única palavra em que as frases diferem, de modo que o contexto
remanescente é o mesmo. O instrumento de medida permanece fixo e apenas o objeto mensurado
varia, e é essa condição que torna a comparação legítima.

A primeira linha exibe o piso da medida. A forma \texttt{destruiu} é a primeira escolha do
modelo, razão pela qual a margem é exatamente nula, o que confere significado ao zero da
escala. Nos casos em que a raiz não é verbo, como no fragmento e na predicação nominal do
corpus, a resposta é \texttt{n/a}, e o reconhecimento da não aplicabilidade integra o emprego
correto do método.

% =====================================================================
%  Método 11 — Embeddings   (dimensão: semântica)
%  Escopo: o que a técnica se propõe a fazer e as fórmulas.
%  A aplicação das fórmulas está na subseção de resultados.
% =====================================================================
\subsection{Embeddings}
\label{met:embeddings}

\indent\indent Todos os métodos anteriores tratam a palavra como símbolo: ela pertence ou não ao léxico, casa
ou não com a regra, possui ou não contagem no corpus, e duas palavras distintas são
simplesmente distintas, sem gradação. Este método se propõe a substituir o símbolo por um
ponto no espaço, obtendo com isso uma noção de proximidade ausente nos demais.

O fundamento é a hipótese distribucional, segundo a qual palavras que ocorrem nos mesmos
contextos possuem significados relacionados, e cuja realização em vetores densos treinados por
predição de contexto é a de~\textcite{mikolov2013}. O vetor sumariza os contextos de ocorrência
da palavra no corpus de treino, e o resultado é uma tabela
$\mathbf{v} : \mathcal{V} \to \mathbb{R}^{d}$ com uma linha por palavra --- aqui a do modelo
\texttt{pt\_core\_news\_md}, com $d = 300$ e 20.000 linhas. Trata-se de consulta, e não de
cálculo, uma vez que o vetor de cada palavra é invariante e está armazenado no modelo. Desta
propriedade decorrem tanto os acertos quanto os limites do método: o vetor não codifica o
significado da palavra, mas os contextos em que ela habitualmente ocorre.

A comparação é feita pelo cosseno, isto é, pelo ângulo entre os dois vetores:
%
\begin{equation}
\cos(\mathbf{a}, \mathbf{b})
= \frac{\mathbf{a} \cdot \mathbf{b}}{\lVert \mathbf{a}\rVert \, \lVert \mathbf{b}\rVert}
\label{eq:cosseno}
\end{equation}
%
O valor situa-se entre $-1$ e $+1$, assumindo $1$ quando os vetores apontam na mesma direção e
$0$ quando são ortogonais. Emprega-se o cosseno, e não a distância euclidiana, porque a norma
$\lVert \mathbf{v} \rVert$ veicula sobretudo frequência, dado que palavras comuns tendem a
vetores de maior comprimento; a divisão pelas duas normas descarta essa magnitude e retém
apenas a direção, onde reside a informação contextual.

A extensão mais simples da medida para a frase é a média dos vetores de suas palavras:
%
\begin{equation}
\vetor{\sent} = \frac{1}{|\sent|} \sum_{t \in \sent} \vetor{t}
\label{eq:vetor-frase}
\end{equation}
%
O somatório percorre um multiconjunto, de modo que qualquer reordenação da frase apresenta
exatamente os mesmos termos e, por consequência, o mesmo vetor. A invariância à ordem é
propriedade da Equação~\ref{eq:vetor-frase}, e não contingência de um exemplo. O método
devolve um número entre $-1$ e $+1$ por par comparado.


% =====================================================================
%  6. Resultados
%  Origem: células 37, 38 e 39 do caderno
%  NOTA: o label sec:empiricos é o alvo dos apontamentos deixados nas
%  subseções de método (m02, m07, m08, m09, m11).
% =====================================================================
\section{Resultados}
\label{sec:empiricos}

Tendo em vista que os melhores resultados foram obtidos com o léxico, com os dois métodos
apoiados no analisador sintático e com o modelo mascarado, os demais não foram incorporados
à cadeia. Os descritores --- índices de complexidade sintática, comprimento de dependência e
embeddings --- não respondem à pergunta da correção, uma vez que medem estilo, custo de
processamento e proximidade temática, e as frases que os separam no corpus são todas
corretas. Os escores globais, SLOR e PLL, exigem um termo de comparação de que a avaliação
de uma frase isolada não dispõe, e o SLOR degenera ainda fora do vocabulário do corpus de
contagens. A surpresa por token, por sua vez, foi excluída por medição direta: seu pico
recaiu sobre o verbo correto da frase de referência.

Aos quatro métodos selecionados incorporaram-se três ajustes identificados nas subseções
anteriores: a exigência de verbo finito na raiz, a comparação de gênero no particípio
passivo e o tratamento do numeral por valor, dado que o numeral não porta traço de número.

Seguindo as etapas lógicas da avaliação, da ortografia à semântica, obtém-se a ordem devida
da cadeia. A verificação ortográfica ocupa a primeira posição por constituir o único
veredito exato e porque uma palavra fora do léxico não permanece contida nessa etapa: ela
deforma a árvore de dependências e, com ela, as duas etapas que a leem. A verificação
estrutural antecede a semântica porque o método do verbo mascarado devolve \texttt{n/a}
quando a raiz não é um verbo; a existência de predicado constitui, portanto, pré-requisito
da etapa seguinte, e o teste da raiz a verifica a custo desprezível.

Cada etapa devolve um conjunto de achados, e o veredito resulta da união
%
\begin{equation}
\mathrm{Loc}(\sent) = E_1(\sent) \cup E_2(\sent) \cup E_3(\sent) \cup E_4(\sent)
\qquad
\mathrm{Loc}(\sent) \neq \varnothing \;\Longrightarrow\; \text{frase reprovada}
\end{equation}
%
O veredito decorre, portanto, da vacuidade de um conjunto: nenhum valor é somado, ponderado
ou confrontado com limiar agregado. A cadeia dispõe de um único parâmetro calibrável, o
$\theta$ da etapa semântica, fixado em 5,0 nats, e ele determina apenas a admissão daquele
achado, sem participar do veredito final. As três primeiras etapas decidem por pertinência a
um conjunto ou por igualdade de valores, sem parâmetro algum.

\subsection*{Percurso da cadeia sobre uma frase}

Aplica-se a seguir a cadeia, etapa a etapa, à frase \emph{Há tres dias o incêndio atingiu
uma áreas de mata natvia}, que reúne três desvios de naturezas distintas. Todas as quatro
etapas são executadas, ainda que a primeira produza achado.

A etapa inicial avalia a pertinência ao léxico, resguardados os nomes próprios e os
numerais:
%
\begin{equation}
E_1(\sent) = \bigl\{\, t \;:\; \ell(t) \notin \lexico,\ \ \mathrm{pos}(t) \neq \texttt{PROPN},\ \ \ell(t) \notin \mathbb{N} \,\bigr\}
\end{equation}
%
O predicado é avaliado sobre uma palavra por vez, contra o conjunto. Das onze palavras, dez
pertencem a $\lexico$ e apenas \texttt{natvia} não, donde $E_1(\sent) =
\{\texttt{natvia}\}$. Cabe registrar a segunda palavra: a forma \texttt{tres}, sem acento,
\textbf{pertence} ao léxico, com $f = 1{.}690$ ocorrências, uma vez que a lista de
frequências foi construída a partir de texto autêntico da web, em que a grafia sem acento
ocorre com frequência suficiente para ser incorporada. A etapa não produz achado nessa
posição, e não por defeito de implementação: trata-se do comportamento previsto pela
fórmula, em que $\lexico$ constitui a definição operacional de erro.

Para o achado detectado, a segunda parte do método computa os candidatos e resolve o
desempate por frequência:
%
\begin{equation}
C(\texttt{natvia}) = \{\texttt{nativa},\ \texttt{naevia},\ \texttt{natia},\ \texttt{latvia}\}
\qquad
\hat{w} = \arg\max_{c \in C} f(c) = \texttt{nativa} \quad (f = 912)
\end{equation}

A segunda etapa verifica a existência de predicado, já sob a exigência de finitude:
%
\begin{equation}
\begin{split}
\mathrm{pred}(\sent) \iff\;
& \bigl[\mathrm{pos}(\raiz) \in \{\texttt{VERB},\texttt{AUX}\} \ \wedge\ \mathrm{VerbForm}(\raiz) = \texttt{Fin}\bigr] \\
\vee\; & \bigl[\exists\, c \in \mathrm{apoio}(\raiz) : \mathrm{VerbForm}(c) = \texttt{Fin}\bigr]
\end{split}
\end{equation}
%
Não há percurso a executar, dado que se inspeciona um único nó. A raiz é \texttt{atingiu},
com $\mathrm{pos} = \texttt{VERB}$ e $\mathrm{VerbForm} = \texttt{Fin}$, e não há auxiliar
nem cópula. Satisfeita a primeira cláusula, a segunda não é avaliada, e $E_2(\sent) =
\varnothing$.

A terceira etapa percorre os arcos e desdobra cada um nos traços exigidos por $T(r)$. Dos
dez arcos da frase, apenas três são admitidos, pois nos demais $T(r) = \varnothing$: as
palavras \texttt{de}, \texttt{mata} e \texttt{áreas} vinculam-se por \texttt{case},
\texttt{nmod} e \texttt{obj}, e nenhuma dessas relações impõe concordância. A
Tabela~\ref{tab:percurso} apresenta as comparações efetuadas.

\begin{table}[H]
\centering
\dimtable{
\begin{tabular}{@{}llccl@{}}
\toprule
\textbf{Arco} $(t,r,h)$ & $\tau$ & $\tracosp{\tau}{t}$ & $\tracosp{\tau}{h}$ & \textbf{Resultado} \\
\midrule
\multirow{2}{*}{(\texttt{o}, det, \texttt{incêndio})}
  & Number & Sing & Sing & iguais \\
  & Gender & Masc & Masc & iguais \\
\midrule
\multirow{2}{*}{(\texttt{incêndio}, nsubj, \texttt{atingiu})}
  & Number & Sing & Sing & iguais \\
  & Person & $\ausente$ & 3 & comparação suprimida \\
\midrule
\multirow{2}{*}{(\texttt{uma}, det, \texttt{áreas})}
  & Number & \textbf{Sing} & \textbf{Plur} & $\in E_3(\sent)$ \\
  & Gender & Fem & Fem & iguais \\
\bottomrule
\end{tabular}}
\caption{As comparações da etapa gramatical sobre os três arcos admitidos.}
\label{tab:percurso}
\end{table}

A quarta linha exemplifica a guarda $\tracosp{\tau}{\cdot} \neq \ausente$: o substantivo
\texttt{incêndio} não porta traço de pessoa, ao passo que o verbo apresenta
$\texttt{Person}=3$, de modo que a comparação é suprimida. O resultado não é de igualdade,
mas de ausência de termos a comparar. Resta, portanto, $E_3(\sent) = \{(\texttt{uma},
\texttt{áreas}, \texttt{Number})\}$.

A quarta etapa inspecciona uma única posição, a da raiz, mascarando-a e comparando a
distribuição obtida com a palavra escrita:
%
\begin{equation}
m = \ln P\bigl(\hat{v} \mid \sent_{\setminus\raiz}\bigr) - \ln P\bigl(v_\raiz \mid \sent_{\setminus\raiz}\bigr)
\end{equation}
%
Mascarada a raiz, o modelo devolve como formas esperadas \texttt{[destruiu, atingiu,
invadiu, atinge, em]}, nas quais o verbo escrito ocupa a segunda colocação, com margem
$m = 0{,}95$. Como $m < \theta = 5{,}0$, tem-se $E_4(\sent) = \varnothing$. Manifesta-se
aqui o único limiar de toda a cadeia: as três etapas anteriores decidem por pertinência a um
conjunto ou por igualdade de valores, ao passo que esta compara um valor real a um ponto de
corte, e $\theta$ constitui escolha do analista.

Reunindo os quatro conjuntos,
%
\begin{equation}
\mathrm{Loc}(\sent) =
\underbrace{\{\texttt{natvia}\}}_{E_1}
\ \cup\ \underbrace{\varnothing}_{E_2}
\ \cup\ \underbrace{\{(\texttt{uma},\,\texttt{áreas})\}}_{E_3}
\ \cup\ \underbrace{\varnothing}_{E_4}
\ \neq\ \varnothing
\end{equation}
%
donde a frase é reprovada, com dois achados, cada qual identificando as palavras envolvidas e
a etapa de origem.

\subsection*{Propagação de erro entre etapas}

A Tabela~\ref{tab:laudo} consolida o laudo dessa frase ao lado do de uma variante que contém
apenas o desvio de concordância.

\begin{table}[H]
\centering
\dimtable{
\begin{tabular}{@{}lll@{}}
\toprule
\textbf{Etapa} & \emph{\ldots uma áreas de mata nativa.} & \emph{Há tres dias \ldots uma áreas \ldots natvia.} \\
\midrule
ortografia & ok \quad --- & \textbf{!!} \quad \texttt{natvia} $\to$ \texttt{nativa} \\
estrutura  & ok \quad predicado presente & ok \quad predicado presente $^\dagger$ \\
gramática  & \textbf{!!} \quad \texttt{uma}(Sing) × \texttt{áreas}(Plur)
           & \textbf{!!} \quad \texttt{uma}(Sing) × \texttt{áreas}(Plur) $^\dagger$ \\
semântica  & ok \quad margem 0,65 & ok \quad margem 0,95 \\
\bottomrule
\end{tabular}}
\caption{Laudo de duas frases. $^\dagger$ marcado como \emph{árvore sob suspeita}.}
\label{tab:laudo}
\end{table}

A decisão de não bloquear tem consequência observável. A palavra \texttt{natvia} não
permanece contida na etapa inicial: ela deforma a árvore que as etapas seguintes leem,
conforme a Tabela~\ref{tab:deformacao}.

\begin{table}[H]
\centering
\dimtable{
\begin{tabular}{@{}lll@{}}
\toprule
\textbf{Arco produzido} & & \textbf{Efeito observado} \\
\midrule
\texttt{(tres, obj, Há)} & & \texttt{tres} analisado como objeto de \texttt{Há} \\
\texttt{(dias, flat:name, tres)} & & \emph{tres dias} lido como nome próprio composto \\
\texttt{(natvia, dep, atingiu)} & & \texttt{dep} é o rótulo de relação não classificável \\
\bottomrule
\end{tabular}}
\caption{Arcos deformados pela palavra fora do léxico.}
\label{tab:deformacao}
\end{table}

Registra-se por isso a marcação \emph{árvore sob suspeita} nas duas etapas que leem a
árvore, marcação que não indica incorreção, mas menor confiabilidade. No caso examinado a
deformação não atingiu o arco relevante, e o achado de concordância é legítimo, o que
constitui, contudo, contingência posicional e não garantia.

Duas lacunas, de naturezas distintas, delimitam o alcance da cadeia. A primeira é a escolha
lexical: \emph{A três dias} atravessa as quatro etapas sem detecção, pois \texttt{a} consta
do léxico, é preposição e não estabelece arco de concordância, e a etapa semântica não
inspeciona essa posição. Trata-se exatamente do caso que as regras escritas à mão detectam,
o que torna sua integração a extensão mais imediata da cadeia. A segunda é o erro factual:
\emph{A área cresceu de 1.284 para 340 hectares} é português correto e falso apenas em
relação à base que originou o texto. Nenhum método de natureza linguística alcança esse
nível, uma vez que a verificação exige confrontar o texto com os campos que o geraram.

% \input{07-pipeline}
% % =====================================================================
%  8. Discussão
%
%  COMPILADO — pontos levantados enquanto as técnicas são percorridas.
%  Nada escrito ainda: cada tópico reúne o achado, sua origem medida e o
%  rótulo onde ele está reportado. A redação vem depois, quando o
%  conjunto estiver fechado.
%
%  A prosa já redigida dos tópicos 1 a 3 está em
%  scratchpad/08-discussao-prosa.tex
% =====================================================================
\section{Discussão}
\label{sec:discussao}


% =====================================================================
%  ESCRITO: Ortografia
% =====================================================================
\subsection{Ortografia}
\label{disc:ortografia}

\indent\indent A dimensão dispõe de um único método, e é ele que devolve o veredito mais forte de todo o
levantamento. A pertinência a $\lexico$ não é estimada nem aproximada: a palavra consta ou não
consta do conjunto, sem gradação e sem margem. É também o único método que, além do
diagnóstico, devolve a correção --- sobre \texttt{natvia}, a Equação~\ref{eq:lex-argmax}
apontou \texttt{nativa} sem ambiguidade, e a subseção~\ref{res:lexico} não registrou falso
positivo sobre a frase correta.

Dois problemas delimitam o que a dimensão alcança, e eles são de naturezas distintas.

O primeiro é de fronteira, e é intrínseco à definição. Em \emph{A três dias} o método não
produziu achado, porque \texttt{a} é palavra do português e pertence a $\lexico$. A Equação~\ref{eq:lex-deteccao} avalia uma palavra por
vez, isoladamente, e \texttt{a}, tomada isoladamente, não apresenta desvio algum.

O segundo problema é de procedência, e afeta o que o método devolve quando acerta. O léxico é
uma lista de frequências, e não um dicionário: o conjunto de candidatos herda o que o corpus
de origem contém. Dos quatro candidatos a uma edição de \texttt{natvia} na
Tabela~\ref{tab:candidatos}, apenas \texttt{nativa} é palavra corrente do português;
\texttt{naevia} e \texttt{natia} não constam de dicionário normativo, e \texttt{latvia} é
forma estrangeira. São candidatos por terem sido atestados no material de origem, e não por
serem palavras do idioma. O desempate segue a mesma lógica, uma vez que a decisão recai
integralmente sobre $f$: \texttt{nativa} vence com 912 ocorrências, então, a qualidade 
do resultado esta diretamente ligada ao corpus utilizado.



% =====================================================================
%  ESCRITO: Gramática
% =====================================================================
\subsection{Gramática}
\label{disc:gramatica}

\indent\indent A dimensão reúne quatro métodos, e eles se dividem em dois grupos pelo tipo de resposta. Os
traços morfossintáticos e as regras escritas à mão devolvem um veredito localizado; o SLOR e o
PLL devolvem um número.

Os dois primeiros acertam onde têm cobertura, sendo assim, dependem das regras bem descritas.
 As regras detectaram a troca em \emph{A três dias} e são o único método do
levantamento a devolver, junto com o achado, o enunciado da norma em português; mas sobre
\emph{os incêndios atingiu} nenhuma das 2.270 regras casou, e $D(\sent) = \varnothing$ é
indistinguível da saída de uma frase correta. É a Equação~\ref{eq:regras-invisivel} verificada
em caso concreto: a cobertura de $R$ é finita e de elaboração manual, e o português dispõe de
249 grupos de regras. Os traços, por sua vez, localizaram o desvio de
\emph{uma áreas} com precisão que nenhum outro método alcança --- a palavra, o termo com que
diverge, o traço violado e o arco que os liga ---, mas deixaram passar \emph{a área inteiro}:
o adjetivo posposto está no masculino e ainda assim $E(\sent) = \varnothing$, porque o
anotador atribuiu $\texttt{Gender} = \texttt{Fem}$ a \texttt{inteiro} e a comparação devolveu
\emph{iguais}. A fórmula está correta e o insumo não está, e é por isso que a
Equação~\ref{eq:tracos-E} compara $\tracosp{}{}$, e não $\tracos{}{}$.

Os dois métodos de escore ordenam, mas não decidem. O PLL ordenou as quatro frases na direção
certa, com degraus amplos, e ainda localizou o desvio, uma vez que $5{,}18$ dos $6{,}00$ nats
de diferença se concentram na posição do verbo. Essa localização, porém, é atribuição do
modelo e não fato da frase: a discordância é \emph{entre} sujeito e verbo, e mascaradas as
formas \texttt{os} e \texttt{incêndios} tornam-se mais prováveis que suas contrapartes no
singular. Nenhum dos dois métodos dispõe de limiar, de modo que os valores são interpretáveis
entre si e nenhum isoladamente.

No SLOR soma-se um problema do estimador. O par de ordem foi separado na direção prevista, mas
a Tabela~\ref{tab:slor} mostra sobre o que a separação repousa: nove das onze parcelas da
frase embaralhada estão no piso $\ln(1-\lambda) = -1{,}204$, e a única contribuição positiva
vem de \texttt{de o}, adjacência entre preposição e artigo que sobreviveu ao embaralhamento
por acidente; do outro lado, 43\% do numerador da frase correta provém de um único bigrama
observado duas vezes. O escore não distingue frase mal construída de frase cujo vocabulário o
corpus de contagens não cobre, e assim nenhum dos quatro métodos decide sozinho sobre
concordância, dois são dependente de regras que podem falhar em casos de cenários 
fora de sua cobertura e dois não dispõem de limiar.



% =====================================================================
%  ESCRITO: Estrutura
% =====================================================================
\subsection{Estrutura}
\label{disc:estrutura}

\indent\indent Os dois métodos da dimensão ocupam posições opostas quanto ao que se pode extrair deles. O
teste da raiz devolve veredito binário e localizado ao custo de inspecionar um único nó; o
comprimento de dependência e a não-projetividade devolvem descritores que não respondem à
pergunta da correção.

O teste da raiz registra uma distinção que nenhum outro método do levantamento alcança. Entre
\emph{o incêndio atingiu} e \emph{o incêndio que atingiu} não há diferença de vocabulário ---
o mesmo verbo está presente nas duas ---, e o que separa oração de fragmento é a posição que
ele ocupa na árvore. Das quatro frases a que foi submetido, o método acertou três, aprovando a
frase de referência pelo primeiro termo da disjunção e a predicação nominal pelo segundo, e
reprovando o fragmento por ambos serem falsos.

O erro é de espécie mais simples que os da dimensão anterior. Em \emph{Atingir uma área de
mata nativa em três dias} a raiz é \texttt{Atingir}, também VERB, e o infinitivo isolado é
aprovado como oração. O traço que separaria os dois casos ---
$\texttt{VerbForm} = \texttt{Inf}$ --- está na árvore, corretamente atribuído, e a
Equação~\ref{eq:raiz-pred} simplesmente não o consulta. Diferentemente do que ocorreu com os
traços morfossintáticos, aqui o insumo está correto e a falha é da definição, portanto
corrigível dentro dela.

Uma ressalva recai sobre a árvore que o método lê. A notação $\raiz(\sent)$ pressupõe
unicidade, e o analisador não a garante: sobre o fragmento devolveu uma floresta, com
\texttt{dias} e \texttt{incêndio} ambos etiquetados \texttt{ROOT}. A implementação retém o
primeiro, e no caso o veredito não depende da escolha, pois os dois candidatos são NOUN --- mas
o acerto se deu sobre uma entrada que a definição não previa.

O segundo método não emite veredito algum, e o corpus deixa isso explícito: as duas frases a
que foi aplicado são corretas e diferem apenas na posição em que as relativas foram
encaixadas. O comprimento médio as separa com margem ampla, de $1{,}42$ para $3{,}58$, com o
maior arco passando de 2 para 9, o que mensura a carga de processamento imposta pela ordem
escolhida. Já a não-projetividade devolveu $X = 0$ nas duas: as desigualdades da
Equação~\ref{eq:dep-cruzam} são estritas, e o encaixe de uma relativa no interior de outra
produz arcos aninhados, não sobrepostos. A medida está correta e sobre este corpus não informa
nada, de modo que da dimensão inteira apenas o teste da raiz é aproveitável para decidir sobre
uma frase.



% =====================================================================
%  ESCRITO: Semântica
% =====================================================================
\subsection{Semântica}
\label{disc:semantica}

\indent\indent A dimensão reúne três métodos, e apenas um deles produziu decisão utilizável.

O verbo mascarado separou o verbo anômalo dos adequados por quase 14 nats, e o zero de sua
escala tem significado, uma vez que \texttt{destruiu} é a primeira escolha do modelo e recebe
margem exatamente nula. Mais relevante que a separação, porém, é a condição que a torna
legítima: a lista de formas esperadas é idêntica nas três frases, porque mascarar a raiz
suprime justamente a única palavra em que elas diferem. O instrumento de medida permanece fixo
enquanto o objeto varia --- propriedade que nenhum outro método do levantamento possui, dado
que no SLOR e na surpresa trocar a palavra altera também as contagens que servem de referência.

Duas ressalvas acompanham o método. A resposta \texttt{n/a} sobre o fragmento e a predicação
nominal é comportamento previsto, e não falha, mas transfere a decisão ao analisador
sintático: as duas frases foram recusadas porque a árvore predita as classificou assim. E
$\mathcal{V}$ não está restrito a formas verbais, como mostra a preposição \texttt{em} entre
as cinco esperadas; neste corpus a folga é inconsequente, pois $\hat{v}$ foi \texttt{destruiu}
nas três frases, mas nada na Equação~\ref{eq:verbo-margem} impede que o melhor preenchimento
seja palavra de outra classe.

A surpresa por token acertou uma das três frases e errou duas, e os três picos têm a mesma
origem. Todos recaem sobre posições de bigrama não observado, onde vale
$\surp_{\mathrm{uni}} + 1{,}204$: na frase de referência, que não contém desvio, o pico caiu
sobre \texttt{atingiu} apenas porque \texttt{incêndio atingiu} não consta do corpus de
contagens. Em \emph{Focos recorrentes apresentaram persistência anômala} o caso é extremo,
pois os cinco bigramas são não observados, todo o perfil é unigrama deslocado de uma constante,
e o pico é a palavra mais rara no valor de teto do estimador. O $\arg\max$ seleciona
raridade, e o acerto sobre \texttt{cantou} é contingência de o desvio ser também uma palavra
rara.

Os embeddings falham por motivo distinto. O bloco de pares sem relação comporta-se como
esperado, situando-se junto a zero, mas o bloco superior não separa o que deveria: o par de
maior proximidade de toda a Tabela~\ref{tab:embeddings} é um par de opostos,
\texttt{seco}~×~\texttt{úmido} com $+0{,}838$, acima de todos os sinônimos e do próprio
\texttt{incêndio}~×~\texttt{fogo}. Não é defeito do modelo, e sim a hipótese distribucional
operando como definida, pois os dois termos ocorrem nos mesmos contextos. Some-se que a
Equação~\ref{eq:vetor-frase} é invariante à ordem --- verificada em $1{,}000000$ sobre o par
de permutação --- e que \texttt{natvia}, fora da tabela de vetores, não recebe medida alguma.

Das três, portanto, apenas o verbo mascarado sobrevive à caracterização, e as duas exclusões se
dão por razões opostas: a surpresa mede raridade em lugar de desvio, e o cosseno mede campo
temático em lugar de sentido.



% =====================================================================
%  ESCRITO: A pipeline (introdução e seleção)
% =====================================================================
\subsection{A pipeline}
\label{disc:pipeline}

\indent\indent Tendo em vista os resultados obtidos, 
observa-se uma abordagem por partição: cada um alcança um tipo 
de desvio e falha em outros. Realizando uma análise segmentada, 
podemos priorizar técnicas em que as falhas não coincidem. 
--- \emph{A três dias} escapa ao léxico e é detectada pelas regras; \emph{os
incêndios atingiu} escapa às regras e é detectada pelo PLL. Encadear os métodos é a resposta
que decorre dessa observação, e o que se apresenta a seguir é uma sugestão de cadeia, não um
resultado do levantamento.

Propõem-se então quatro etapas, uma por dimensão, cada qual ocupada pelo método que a discussão
anterior identificou como o único da sua dimensão a produzir veredito localizado sem parâmetro
a calibrar:

\begin{center}
\small
ortografia $\;\rightarrow\;$ estrutura $\;\rightarrow\;$ gramática $\;\rightarrow\;$ semântica \\[2pt]
\emph{léxico} \qquad \emph{teste da raiz} \qquad \emph{traços} \qquad \emph{verbo mascarado}
\end{center}

A primeira etapa é o léxico porque é o único veredito do levantamento que não depende de
estimativa: a palavra consta ou não consta de $\lexico$. É também o único método que devolve a
correção junto com o diagnóstico. A ressalva da subseção~\ref{disc:ortografia} permanece --- a
qualidade do que ele devolve está ligada ao corpus que originou a lista ---, mas ela afeta a
correção sugerida, não a detecção.

A segunda é o teste da raiz porque, das duas técnicas estruturais, é a única que emite
veredito: o comprimento de dependência mede carga de processamento sobre frases corretas, e a
não-projetividade não separou nada. Pesa a favor dele que sua única falha, a aprovação do
infinitivo isolado, seja corrigível dentro da própria definição, bastando que a
Equação~\ref{eq:raiz-pred} passe a exigir $\texttt{VerbForm} = \texttt{Fin}$ na raiz. É a
correção que a cadeia pressupõe.

A terceira são os traços morfossintáticos, e não as regras escritas à mão, embora as duas
detectem desvios reais. A razão é a forma da cobertura: $T(r)$ é uma condição sobre classes de
arcos e vale para toda ocorrência da classe, ao passo que $R$ é um conjunto finito de padrões
individuais, cuja ausência de item torna o desvio invisível. Os traços também localizam com a
precisão mais fina do levantamento. As regras permanecem, contudo, a extensão mais imediata da
cadeia, uma vez que detectam a escolha lexical, que é o desvio que nenhuma das quatro etapas
alcança.

A quarta é o verbo mascarado por ser o único método semântico que sobreviveu à caracterização,
e porque é o único de todo o levantamento cujo instrumento de medida permanece fixo enquanto o
objeto varia.

A ordem das quatro não é escolha de projeto. A ortografia vem primeiro porque uma palavra fora
do léxico não permanece contida nessa etapa: ela deforma a árvore de dependências que as três
etapas seguintes leem. A estrutura precede a gramática porque a existência de predicado é
condição para que os arcos de concordância estejam corretamente atribuídos, e o teste da raiz a
verifica a custo desprezível. E a semântica encerra a cadeia porque o verbo mascarado devolve
\texttt{n/a} quando a raiz não é verbo, de modo que depende da verificação estrutural. Cabe
registrar que essa dependência não é estrita: uma predicação nominal passa pela segunda etapa e
ainda assim recebe \texttt{n/a} na quarta, pois ter predicado e ter verbo na raiz não são a
mesma condição.


Aplicada a cadeia, cada etapa devolve um achado ou nada, e a frase é reprovada quando alguma
delas produz achado. As quatro são executadas sempre, ainda que a primeira já reprove: a
decisão de não bloquear é o que torna o laudo legível, pois expõe todos os desvios de uma vez
em lugar de um por passagem. A Tabela~\ref{tab:laudo} apresenta o laudo de duas frases do
corpus.

\begin{table}[H]
\centering
\dimtable{
\begin{tabular}{@{}llll@{}}
\toprule
& \textbf{Etapa} & & \textbf{Achado} \\
\midrule
\multicolumn{4}{@{}l}{\emph{Há três dias o incêndio atingiu uma áreas de mata nativa.}} \\
\midrule
& ortografia & ok & --- \\
& estrutura  & ok & predicado presente \\
& gramática  & \textbf{!!} & Number: \texttt{uma}(Sing) × \texttt{áreas}(Plur) [det] \\
& semântica  & ok & \texttt{atingiu}, margem 0,65 \\
\midrule
\midrule
\multicolumn{4}{@{}l}{\emph{Há tres dias o incêndio atingiu uma áreas de mata natvia.}} \\
\midrule
& ortografia & \textbf{!!} & \texttt{natvia} $\to$ \texttt{nativa} \\
& estrutura  & ok & predicado presente \quad$^\dagger$ \\
& gramática  & \textbf{!!} & Number: \texttt{uma}(Sing) × \texttt{áreas}(Plur) [det] \quad$^\dagger$ \\
& semântica  & ok & \texttt{atingiu}, margem 0,95 \quad$^\dagger$ \\
\bottomrule
\end{tabular}}
\caption{Laudo da cadeia sobre duas frases. $^\dagger$ etapa marcada como \emph{árvore sob
suspeita}, por ler uma árvore construída sobre palavra fora do léxico.}
\label{tab:laudo}
\end{table}

A primeira frase contém um único desvio, e a cadeia o localiza na etapa que lhe corresponde. A
segunda acrescenta uma palavra fora do léxico, e é ela que expõe o custo do encadeamento: a
forma \texttt{natvia} não permanece contida na primeira etapa, pois deforma a árvore de
dependências que as três seguintes leem. Daí a marcação \emph{árvore sob suspeita}, que não
indica incorreção, mas confiabilidade menor --- e que é a resposta prática ao limite enunciado
na subseção~\ref{disc:gramatica}, o de que uma cadeia de etapas que leiam a mesma árvore
predita herda os erros dessa predição.

No caso examinado a deformação não atingiu o arco relevante, e o achado de concordância
permanece legítimo nas duas frases. Trata-se, contudo, de contingência posicional, e não de
garantia: a marcação existe justamente porque a cadeia não tem como distinguir os dois casos.

Cabe observar, por fim, que a quarta etapa não produz achado em nenhuma das duas frases, com
margens de $0{,}65$ e $0{,}95$. O ganho dela não se manifesta neste par, cujos desvios são de
ortografia e de concordância; ela se destina à anomalia de seleção verbal, para a qual a margem
medida na subseção~\ref{res:verbo} alcança $13{,}73$.



% =====================================================================
%  TÓPICO 1 — [ESCRITO em \subsection{Ortografia}]
%  origem: res:lexico, tab:candidatos
% ---------------------------------------------------------------------
%  - O léxico é lista de frequências, não dicionário. O conjunto de
%    candidatos herda o que o corpus contém: dos quatro candidatos de
%    'natvia', só 'nativa' é palavra corrente do português — 'naevia' e
%    'natia' não constam de dicionário normativo, 'latvia' é estrangeira.
%  - O desempate reproduz a distribuição do corpus: 'nativa' vence com
%    f = 912, margem do material de origem e não propriedade linguística.
%  - Trocar a lista altera quem entra em C(w) E quem vence o argmax, sem
%    que uma linha da definição mude.
%  - Consequência: a qualidade da correção é herdada, não construída. O
%    veredito de pertinência a D — o único exato entre os métodos — é
%    exato apenas relativamente ao conjunto que se adotou.


% =====================================================================
%  TÓPICO 2 — Três falhas, três causas
%  origem: res:lexico, res:regras, e o caso 'a área inteiro' (célula 9)
% ---------------------------------------------------------------------
%  Cada método deixou passar um desvio, e as causas estão em níveis
%  diferentes. É isso que justifica encadear em vez de aperfeiçoar um.
%
%  (a) DEFINIÇÃO — 'A três dias', no léxico.
%      'a' pertence a D; eq:lex-deteccao se esgota nesse fato. Não há
%      defeito a corrigir: predicado que avalia uma palavra isolada não
%      distingue forma existente de forma adequada.
%
%  (b) COBERTURA DO RECURSO — 'os incêndios atingiu', nas regras.
%      Nenhuma das 2.270 regras casou; D(s) = vazio. eq:regras-invisivel
%      previa exatamente isso. Correção disponível: escrever mais uma
%      regra — trabalho manual que não generaliza.
%
%  (c) INSUMO — 'a área inteiro', nos traços.  <-- a mais difícil
%      Adjetivo posposto no masculino, núcleo no feminino, e ainda assim
%      E(s) = vazio. O arco amod foi admitido por T(r), a comparação de
%      Gender foi efetuada e devolveu 'iguais': o anotador atribuíra
%      Gender=Fem a 'inteiro'. eq:tracos-E compara mu-chapéu, não mu.
%      A formulação está correta; nenhum ajuste dela recupera o caso.
%      É a promessa que met:tracos deixa em aberto e que aponta para cá.


% =====================================================================
%  TÓPICO 3 — O que as falhas justificam, e o que elas não justificam
%  origem: res:lexico + res:regras (par medido de complementaridade)
% ---------------------------------------------------------------------
%  - Nenhuma das três falhas se corrige por dentro do método que a
%    produziu, e elas não coincidem: 'A três dias' escapa ao léxico e é
%    detectada pelas regras. É o argumento direto para o encadeamento —
%    não porque um método seja melhor, mas porque operam sobre objetos
%    distintos e erram em lugares distintos.
%  - LIMITE, a enunciar junto: a falha (c) é a que o encadeamento NÃO
%    resolve. Etapas que leiam a mesma árvore predita herdam o erro dessa
%    predição. O encadeamento resolve complementaridade de cobertura, não
%    qualidade do insumo.
%  - Daí a necessidade de uma etapa registrar quando há motivo para
%    suspeitar da árvore, em vez de tratar todo veredito como igualmente
%    confiável. (Prepara a marcação 'árvore sob suspeita' da pipeline.)


% =====================================================================
%  TÓPICO 4 — O SLOR não ordena boa formação
%  origem: res:slor, tab:slor
% ---------------------------------------------------------------------
%  - PAR MÍNIMO SEM SEPARAÇÃO: 'Os satélites registraram/registrou 47
%    focos' recebem o MESMO escore, -0,072, até a 3a casa. Nenhum dos dois
%    bigramas foi observado; ambas as posições recuam para a unigrama.
%  - FRASE CORRETA NO PISO: 'Focos recorrentes apresentaram persistência
%    anômala' = -1,204 = ln(1-lambda), exatamente eq:piso-slor. Nenhum de
%    seus 5 bigramas observado. Mede ausência de vocabulário no corpus de
%    contagens, não construção.
%  - ORDENAÇÃO INVERTIDA: a frase correta pontua ABAIXO da embaralhada
%    (-1,000) e ABAIXO da frase sem sentido de Chomsky (-0,939).
%  - TESE: a raridade que o método foi construído para descontar reaparece
%    pelo piso do estimador. Liga ao TÓPICO 1 — de novo o corpus de
%    contagens determinando o veredito, agora num método de escore.
%  - Consequência para a cadeia: escore global sem termo de comparação não
%    decide sobre frase isolada. Reforça o veredito por vacuidade de
%    conjunto em vez de escore agregado.


% =====================================================================
%  A ACRESCENTAR conforme as técnicas forem percorridas:
%  PLL, raiz, dependência, surpresa, verbo, embeddings
% =====================================================================

% % =====================================================================
%  9. Conclusão
% =====================================================================
\section{Conclusão}
\label{sec:conclusao}

\indent\indent Este trabalho descreveu dez métodos de avaliação de frases isoladas em português sob um
formalismo único e os submeteu ao mesmo conjunto controlado de variantes, distribuídos por
quatro dimensões: ortografia, gramática, estrutura e semântica.

A caracterização não produziu uma ordenação. Nenhum dos dez domina os demais, e as falhas
observadas não se corrigem umas às outras porque nascem em níveis distintos: umas na definição
do próprio método, outras na cobertura finita do recurso que ele consulta, outras ainda no
insumo que um modelo anterior lhe entrega.

Sendo assim decorre que combinar os métodos não é conveniência de implementação, mas a resposta que o
levantamento indica. Como as coberturas são complementares e as falhas não coincidem, o que
um método deixa passar é frequentemente o que outro alcança --- e a cadeia proposta reúne
quatro etapas, uma por dimensão, decidindo pela existência de achados localizados em lugar de
escore agregado. Cada item do laudo conserva a etapa que o produziu e as palavras envolvidas, e
um único limiar figura em toda a cadeia.

Os limites são de duas ordens, e ambas recomendam cautela na leitura. O corpus é controlado e
pequeno, derivado de uma única frase de referência, e cada método foi aplicado às variantes que
exercitam a propriedade que ele mede, de modo que a complementaridade permanece argumento de
construção mais do que medida. E os instrumentos são únicos --- um analisador sintático, um
corpus de contagens, um modelo mascarado ---, nenhum treinado ou ajustado aqui, o que significa
que parte do que se observou pertence a eles, e não às definições.

Desses limites decorrem os trabalhos futuros. O primeiro é completar a aplicação cruzada,
submetendo cada método a todas as variantes, para que a complementaridade deixe de ser
argumento e passe a ser medida. O segundo é ampliar o corpus, tanto em número de frases de
referência quanto em tipos de desvio, e submetê-lo a julgamentos humanos de aceitabilidade,
de que este trabalho não dispôs. O terceiro é variar os instrumentos, repetindo as medições sob
outro analisador, outro corpus de contagens e outro modelo mascarado, o que permitiria separar
o que pertence às definições do que pertence às ferramentas. O quarto é incorporar as regras
escritas à mão à cadeia, uma vez que alcançam a escolha lexical, desvio que as quatro etapas
atuais atravessam sem detectar. Resta, por fim, o erro factual, que é português correto e falso
apenas em relação à base que originou o texto: sua verificação exige confrontar o enunciado com
os campos que o geraram, e devolve o problema ao sistema gerador.


% ---------------------------------------------------------------------
\printbibliography[title={Referências}]

% ---------------------------------------------------------------------
% \appendix
% \input{A-codigo}
% \input{B-tabelas-derivacao}
% \input{C-protocolo}

\end{document}
